% !TeX root = ../main.tex
\chapter{结论与展望}\label{chap:conclusion}

\section{研究工作总结}

云原生环境下的权限提升漏洞从披露到修复往往面临显著的补丁窗口期。在此阶段，生产环境迫切需要可快速部署的临时防控方案以压制攻击路径、降低风险暴露面。然而，现有安全工具与方法多侧重通用规则检测或静态配置核查，难以针对特定 $\mathrm{CVE}$ 的利用前置条件与攻击链路给出精准且可落地的处置建议。面向这一现实需求，本文提出基于大语言模型的多智能体协同漏洞策略化防控框架，将漏洞语义理解、安全知识检索、策略生成与质量评估有机融合，实现从 $\mathrm{CVE}$ 公告到可执行防控策略的全流程自动化。

围绕"证据可追溯、结构化产物、质量闭环"的核心设计理念，本文在理论建模、方法体系与工程实践三个层面均取得实质性进展。主要贡献与成果概括如下：

\subsection{理论建模与形式化框架}

\noindent\hangindent=2em\hangafter=1 \textbf{提出监控策略与防护策略两分类体系及形式化表示}。基于云原生安全防控机制的作用时机与机理差异，本文将安全策略划分为监控策略（$\Pi_{\mathrm{mon}}$）与防护策略（$\Pi_{\mathrm{prot}}$）两个互补类别：监控策略侧重运行时异常行为的持续观测与取证，防护策略侧重事前阻断与攻击面收敛。每个策略 $\pi$ 均表达为四元组 $\langle \text{type},\; \text{layer},\; \text{method},\; \text{actions} \rangle$，其中 $\text{type}$ 指明策略类型，$\text{layer}$ 描述防控技术层面，$\text{method}$ 刻画防控思路分类，$\text{actions}$ 为可执行的防控操作序列。该形式化定义不仅为生成式模型提供了结构化产出约束，也为后续策略验证与工程落地建立了统一的接口规范。

\noindent\hangindent=2em\hangafter=1 \textbf{建立基于排名的三维质量评价机制}。针对候选策略质量刻画与选择问题，本文定义三维质量向量 $q(\pi) = (E(\pi),\; I(\pi),\; D(\pi))$，分别衡量策略的有效性、无干扰性与可部署性。在此基础上，引入基于排名的评价机制：对每一维度构建严格排序 $R_{\delta}(\Pi_v)$，并通过秩函数 $r_{\delta}$ 与归一化得分函数 $s_{\delta}(\pi) = \frac{n - r_{\delta}(\pi)}{n - 1}$ 将相对次序转化为可计算的数值表示。该机制采用相对刻度而非绝对评分，有效避免了跨样本对比时的尺度不一致问题，更适合在不同漏洞场景中进行可比性分析与鲁棒决策。

\noindent\hangindent=2em\hangafter=1 \textbf{形式化定义多评审主体共识排序与聚合模型}。为抑制单一评审视角的偏差与不一致性，本文提出多主体交叉评审与加权聚合机制。每个评审主体 $h\in\mathcal{H}$ 对候选集合 $\Pi_v$ 给出独立的维度排名 $r_{\delta}^{(h)}(\pi)$，随后通过加权平均计算共识得分 $\hat{S}_{\delta}(\pi)=\sum_{h\in\mathcal{H}} w_h\cdot s_{\delta}^{(h)}(\pi)$，最终基于多维度共识得分实施综合排序。该模型通过结构化的数学表达保证了评价过程的可审计性与可复现性。

\subsection{方法体系与技术创新}

\noindent\hangindent=2em\hangafter=1 \textbf{构建面向防控的云原生权限提升漏洞分类框架与安全知识库}。针对 $\mathrm{CWE}$ 标签非正交、抽象层次不一致的问题，本文提出两级成因分类框架：一级分类刻画成因大类（安全逻辑缺陷、配置缺陷、敏感信息泄露、代码注入），二级分类细化为更贴近控制动作选择的类型（如授权验证不完善、冗余权限、输入验证不足等）。该框架不仅为策略生成提供了稳定且可解释的风险语义约束，也有助于将"漏洞语义$\rightarrow$控制域$\rightarrow$可执行动作序列"对齐，提高输出策略的针对性与可操作性。在知识库构建层面，本文从开源代码仓库、$\mathrm{CVE}$ 公告、$\mathrm{Issue}$ 与 $\mathrm{Pull\ Request}$ 等多源渠道汇聚安全证据，通过分块、去重与向量化索引构建可检索的上下文知识库，确保证据来源可标注、可回溯，为生成式推理提供结构化先验约束。

\noindent\hangindent=2em\hangafter=1 \textbf{设计多智能体协同的策略生成流程与迭代优化机制}。本文提出"知识检索$\rightarrow$威胁分析$\rightarrow$策略生成$\rightarrow$评审反馈"四阶段生成流水线，采用"$\mathrm{RAG}$+$\mathrm{CoT}$+$\mathrm{Feedback}$"技术组合实现可控生成。具体而言，安全知识智能体负责从知识库检索相关证据并整理上下文；漏洞分析智能体基于 $\mathrm{CoT}$ 推理输出结构化威胁分析报告；策略生成智能体分别产出监控策略与防护策略草案；策略评审智能体与策略优化智能体通过多轮反馈引导候选策略迭代修订。该分阶段设计将复杂的端到端生成任务分解为可解释的推理步骤，显著降低了输出噪声与逻辑不一致性，使生成策略更贴近工程落地需求。

\noindent\hangindent=2em\hangafter=1 \textbf{实现静态评估与动态验证相结合的两阶段质量控制}。在静态策略评估阶段，多个评价智能体基于三维质量向量对候选策略进行交叉评审，输出维度排名与共识排序；在动态策略验证阶段，针对监控策略通过沙箱环境注入攻击行为并评估监控有效性，针对防护策略通过模拟漏洞利用路径验证阻断能力与业务兼容性。两阶段验证结果共同构成策略筛选与排序的依据，既保证了策略质量的多视角校验，也为最终交付提供了可审计的质量证明。

\subsection{工程实践与实验验证}

\noindent\hangindent=2em\hangafter=1 \textbf{实现覆盖知识库构建、流程编排与策略管理的平台化工具体系}。为支撑上述方法体系的工程化落地，本文设计并实现了漏洞安全策略生成与管理平台，涵盖安全知识库构建、工作流管理、模型对接、人工交互与策略部署管理五项核心能力。该平台通过模块化架构实现知识库维护与策略生成流程的解耦，支持多种大语言模型的灵活对接与评价主体的动态配置，并提供可视化界面用于人工复核与策略迭代优化，从而将智能防控能力转化为可持续运营的安全服务。

\noindent\hangindent=2em\hangafter=1 \textbf{在真实 $\mathrm{CVE}$ 样本上开展系统性实验评估并验证方法有效性}。本文对 2020 年至 2025 年间的 55 个云原生权限提升漏洞样本进行收集、标注与分类，构建涵盖 $\mathrm{Kubernetes}$、$\mathrm{Cilium}$、$\mathrm{etcd}$、$\mathrm{Harbor}$ 等核心组件与第三方应用的实验数据集。实验结果表明：（1）本框架生成的策略在有效性、无干扰性与可部署性三个维度均显著优于基线方法，其中有效性评分提升 32.5\%，可部署性评分提升 41.2\%；（2）知识库检索、$\mathrm{CoT}$ 推理与反馈优化三个模块均对最终交付质量产生显著正向贡献，消融实验证实各模块缺失均会导致策略质量明显下降；（3）多评审主体机制能够有效抑制单一模型的偏差，共识排序结果较单一评审的稳定性提升 23.7\%；（4）案例分析进一步展示了框架在 $\mathrm{CVE}$-2023-30512、$\mathrm{CVE}$-2023-3955 等典型漏洞上的实际应用效果，生成的监控策略与防护策略均可在 $\mathrm{Falco}$、$\mathrm{Kyverno}$、$\mathrm{NetworkPolicy}$ 等工具上快速部署并发挥预期作用。

\subsection{研究意义与价值}

本文工作在理论、方法与实践三个层面均具有重要意义：

\textbf{理论层面}，本文为云原生漏洞防控提供了形式化建模框架与质量评价机制，将"漏洞语义$\rightarrow$安全策略$\rightarrow$质量度量"的全流程纳入统一的数学表达体系，为后续研究提供了可扩展的理论基础。基于排名的评价机制与多主体共识聚合模型在跨样本可比性与鲁棒决策方面表现出明显优势，可为其他安全策略生成场景提供方法论参考。

\textbf{方法层面}，本文提出的多智能体协同生成流程与两阶段质量控制机制，有效解决了大语言模型在安全策略生成中面临的幻觉、噪声与可控性不足等问题。通过将生成任务分解为可解释的推理步骤，结合结构化约束与多轮反馈优化，本框架在保持生成灵活性的同时显著提升了输出质量与可落地性，为生成式$\mathrm{AI}$在安全工程中的应用探索了可行路径。

\textbf{实践层面}，本文工作直接服务于云原生环境下权限提升漏洞的补丁窗口期防控需求，生成的监控策略与防护策略可在 $\mathrm{Kubernetes}$ 生产环境中快速部署，有效压制攻击路径并降低风险暴露面。实验结果与案例分析表明，本框架较现有静态工具在策略针对性、可部署性与业务兼容性方面均具有明显优势，可为企业安全团队在补丁尚未就绪或升级受阻阶段提供及时、有效的临时防控方案，具有重要的工程应用价值。

综上所述，本文围绕云原生权限提升漏洞的补丁窗口期防控问题，提出了完整的理论框架、方法体系与工程方案，并通过系统性实验验证了方法的有效性与实用性。研究成果不仅为云原生安全治理提供了新的技术手段，也为生成式$\mathrm{AI}$在安全工程中的应用积累了宝贵经验。

\section{未来研究方向}

尽管本文工作在云原生权限提升漏洞的补丁窗口期防控方面取得了积极进展，但仍存在进一步提升与扩展的空间。未来研究可从以下三个方向展开：

\subsection{高危漏洞的自动化监控与实时更新}

当前知识库采用离线批量构建方式，无法及时响应新披露的权限提升漏洞。未来可构建面向 $\mathrm{GitHub}$ 仓库与 $\mathrm{NVD}$ 漏洞库的持续监控机制，通过以下技术手段实现漏洞情报的自动化获取与更新：
\begin{itemize}
    \item \textbf{多源漏洞监控}：对 $\mathrm{GitHub\ Security\ Advisory}$、$\mathrm{NVD}$、$\mathrm{MITRE\ CVE}$ 等公开数据源进行实时爬取与订阅，结合关键词过滤（如 Privilege Escalation、Authorization Bypass、Container Escape 等）与 $\mathrm{CVSS}$ 评分筛选，自动识别云原生环境下的高危权限提升漏洞；
    \item \textbf{增量知识库更新}：当检测到新漏洞披露时，自动触发知识库构建流程，从相关代码仓库、$\mathrm{Issue}$ 讨论与补丁提交中提取上下文证据，并更新向量索引与分类标签，确保策略生成所依赖的安全知识保持最新状态；
    \item \textbf{变更通知与策略重评估}：对已生成策略的漏洞，当检测到补丁发布、受影响版本范围更新或新的利用路径披露时，自动触发策略重评估流程，并向运维团队推送变更通知，辅助其判断是否需要调整现有防控措施。
\end{itemize}
该机制可显著缩短从漏洞披露到策略交付的响应时间，使防控体系能够更快速地适应不断演变的威胁态势。

\subsection{漏洞利用的自动化复现与验证}

当前策略验证主要依赖人工模拟攻击或基于文档描述的静态推理，缺乏真实漏洞利用环境下的自动化验证能力。未来可构建漏洞自动化复现框架，通过以下技术路径提升策略验证的准确性与可信度：
\begin{itemize}
    \item \textbf{漏洞环境自动化搭建}：基于漏洞披露信息与受影响组件版本，自动构建包含漏洞的 $\mathrm{Kubernetes}$ 测试集群（如通过 $\mathrm{Kind}$、$\mathrm{Minikube}$ 或容器化部署），并配置与生产环境一致的网络拓扑、$\mathrm{RBAC}$ 策略与工作负载；
    \item \textbf{利用脚本生成与执行}：结合漏洞 $\mathrm{PoC}$（Proof of Concept）代码、公开 $\mathrm{Exploit}$ 脚本与漏洞分析报告，自动生成或适配利用代码，并在隔离环境中执行攻击以验证漏洞可复现性；
    \item \textbf{策略有效性量化评估}：在已复现漏洞的环境中部署生成的监控策略与防护策略，通过重放攻击行为评估策略的阻断成功率、监控覆盖率与误报率，并生成包含攻击日志、策略触发记录与业务影响分析的可视化验证报告。
\end{itemize}
该能力可为策略质量评价提供更客观的证据支撑，同时也有助于识别策略设计中的盲点与不足，指导迭代优化。

\subsection{研究范围向云原生其他漏洞类型扩展}

本文聚焦于权限提升漏洞的防控，未来可将方法体系扩展至云原生环境中的其他高危漏洞类型，以构建更全面的安全防护能力：
\begin{itemize}
    \item \textbf{拒绝服务（DoS）漏洞}：面向资源耗尽、崩溃重启等拒绝服务攻击，生成基于资源配额限制、$\mathrm{Pod\ Disruption\ Budget}$ 与速率控制的防护策略，并结合监控告警机制实现异常流量的快速检测与隔离；
    \item \textbf{敏感信息泄露漏洞}：针对 $\mathrm{Secret}$ 泄露、日志敏感信息暴露、$\mathrm{API}$ 端点未授权访问等场景，生成基于最小权限原则、日志脱敏与网络访问控制的防护策略，并通过运行时监控检测异常数据传输行为；
    \item \textbf{供应链与镜像安全漏洞}：面向恶意镜像注入、依赖组件漏洞等供应链风险，生成基于镜像签名验证、$\mathrm{SBOM}$ 审计与准入控制的防护策略，并结合漏洞扫描工具实现持续化的安全基线核查。
\end{itemize}
通过扩展漏洞分类框架、丰富策略生成模板与优化质量评价维度，本框架可为更广泛的云原生安全场景提供智能化防控支撑。

综上所述，未来研究将围绕"实时响应、自动化验证、全面覆盖"三个核心目标展开，推动云原生安全防控从被动应对向主动防御、从单点防护向体系化防御的持续演进。
